\documentclass[12pt, a4paper]{article}
\usepackage[brazil]{babel}
\usepackage[utf8]{inputenc}
\usepackage[T1]{fontenc}
\usepackage{amsmath, amssymb, amsthm}
\usepackage{geometry}
\usepackage{booktabs}
\usepackage{float}
\usepackage{graphicx}
\usepackage{tikz}
\usetikzlibrary{calc, arrows.meta, positioning, shapes, decorations.pathmorphing, decorations.markings}
\usepackage{hyperref}
\usepackage{cleveref}
\usepackage{physics}
\usepackage{cancel}

\geometry{a4paper, margin=2.5cm}

% Definindo ambientes matemáticos
\newtheorem{definition}{Definição}
\newtheorem{theorem}{Teorema}
\newtheorem{lemma}{Lema}
\newtheorem{corollary}{Corolário}
\newtheorem{example}{Exemplo}
\newtheorem{remark}{Observação}
\newtheorem{axiom}{Axioma}
\newtheorem{proposition}{Proposição}

\title{Integrais de Linha na Dinâmica Emocional: \\
       Geometria, Topologia e Teoremas Fundamentais}
\author{Do Cálculo Vetorial à Topologia do Espaço Emocional}
\date{Versão 1.0 - Análise Geométrica de Trajetórias Emocionais}

\begin{document}

\maketitle

\begin{abstract}
Este artigo desenvolve uma teoria geométrica das trajetórias emocionais através do conceito de integral de linha. Definimos a integral de caminho do custo existencial ao longo de trajetórias emocionais e exploramos seus significados psicológicos. Apresentamos teoremas fundamentais (Green, Stokes, Gauss) adaptados ao contexto emocional, demonstrando como a topologia do espaço emocional afeta a dinâmica de evolução. Estabelecemos conexões entre homotopia emocional, campos conservativos e a existência de traumas como singularidades topológicas.
\end{abstract}

\tableofcontents

\section{Introdução: Geometria das Experiências}

A vida emocional pode ser vista como uma trajetória contínua no espaço de estados emocionais $\mathcal{E}$. Cada experiência, cada interação, desenha um caminho neste espaço. A integral de linha nos permite quantificar o "custo total" ou "aprendizado acumulado" ao longo dessas trajetórias, revelando estruturas profundas da dinâmica emocional.

\begin{center}
\begin{tikzpicture}[scale=1.2]
  % Espaço emocional
  \draw[gray, very thin] (-1,-1) grid (4,3);
  \draw[->, thick] (-1,0) -- (4.5,0) node[right] {Alegria-Tristeza};
  \draw[->, thick] (0,-1) -- (0,3.5) node[above] {Calma-Agitação};
  
  % Caminho emocional
  \draw[red, thick, ->, decoration={markings, mark=at position 0.3 with {\arrow{>}}, 
        mark=at position 0.7 with {\arrow{>}}}, postaction={decorate}] 
     plot[smooth, tension=0.7] coordinates {(0.5,0.5) (1,2) (2.5,2.5) (3.5,1)};
  
  % Pontos ao longo do caminho
  \fill[red] (0.5,0.5) circle (2pt) node[below left] {$x_0$};
  \fill[red] (3.5,1) circle (2pt) node[above right] {$x_1$};
  
  % Elemento infinitesimal
  \draw[blue, thick, ->] (2,2.54) -- node[above, sloped] {$dx$} (2.3,2.74);
  \fill[blue] (2,2.54) circle (1pt);
  
  % Campo vetorial (gradiente de V)
  \foreach \x in {0.5,1.5,2.5,3.5}
    \foreach \y in {0.5,1.5,2.5}
      \draw[green!50!black, ->] (\x,\y) -- +(-0.2,-0.1);
  
  \node at (1,3) {$\nabla V(x)$};
  \node at (2, -0.5) {Trajetória emocional $\gamma(t)$};
\end{tikzpicture}
\end{center}

\section{Integral de Linha do Custo Existencial}

\subsection{Definição Formal}

\begin{definition}[Trajetória Emocional]
Uma \textbf{trajetória emocional} é uma curva suave por partes:
\[
\gamma: [a,b] \to \mathcal{E}, \quad t \mapsto \gamma(t)
\]
onde $\gamma(a) = x_0$ é o estado inicial e $\gamma(b) = x_1$ é o estado final.
\end{definition}

\begin{definition}[Integral de Linha do Custo]
Dada uma função $V: \mathcal{E} \to \mathbb{R}$ (custo existencial) e uma trajetória $\gamma$, a \textbf{integral de linha} de $V$ ao longo de $\gamma$ é:
\[
\int_\gamma V \, ds = \int_a^b V(\gamma(t)) \|\gamma'(t)\| \, dt
\]
onde $\|\gamma'(t)\|$ é a velocidade emocional instantânea.
\end{definition}

\begin{remark}[Interpretação Psicológica]
Se $V(x)$ mede o "sofrimento" ou "custo" do estado $x$, então $\int_\gamma V \, ds$ representa o \textbf{sofrimento acumulado} ao longo da trajetória $\gamma$. Não depende apenas dos estados visitados, mas também da velocidade com que se transita entre eles.
\end{remark}

\subsection{Exemplo: Caminho Direto vs. Caminho Tortuoso}

\begin{example}[Dois caminhos entre os mesmos estados]
Considere $V(x,y) = x^2 + y^2$ (custo cresce com a intensidade emocional) e dois caminhos de $(-1,0)$ a $(1,0)$:

\begin{enumerate}
\item \textbf{Caminho direto}: $\gamma_1(t) = (2t-1, 0)$, $t \in [0,1]$
\[
\int_{\gamma_1} V \, ds = \int_0^1 [(2t-1)^2 + 0^2] \cdot 2 \, dt = 2\int_0^1 (4t^2 - 4t + 1) \, dt = \frac{2}{3}
\]

\item \textbf{Caminho semicircular}: $\gamma_2(t) = (\cos(\pi t), \sin(\pi t))$, $t \in [0,1]$
\[
\int_{\gamma_2} V \, ds = \int_0^1 [\cos^2(\pi t) + \sin^2(\pi t)] \cdot \pi \, dt = \int_0^1 \pi \, dt = \pi
\end{enumerate}

O caminho direto (menor intensidade) acumula menos custo que o caminho intenso (semicírculo).
\end{example}

\section{Integral de Linha de Campos Vetoriais}

\subsection{Campos de Força Emocional}

\begin{definition}[Campo Vetorial Emocional]
Um \textbf{campo vetorial} em $\mathcal{E}$ é uma função $F: \mathcal{E} \to \mathbb{R}^n$ que associa a cada estado emocional $x$ um vetor $F(x)$ representando uma "força" ou "tendência" emocional.
\end{definition}

\begin{example}[Campo de Retorno aos Valores]
O campo associado à meditação é:
\[
F_{\text{med}}(x) = -\varepsilon (x - \Pi_{\mathcal{V}}(x))
\]
que aponta na direção dos valores $\mathcal{V}$.
\end{example}

\begin{definition}[Integral de Linha de um Campo]
Dado um campo vetorial $F$ e uma curva $\gamma$, a \textbf{integral de linha} de $F$ ao longo de $\gamma$ é:
\[
\int_\gamma F \cdot d\vec{r} = \int_a^b F(\gamma(t)) \cdot \gamma'(t) \, dt
\]
onde $\cdot$ é o produto escalar.
\end{definition}

\begin{remark}[Interpretação Física]
Esta integral mede o \textbf{trabalho} realizado pelo campo $F$ ao longo do caminho $\gamma$. No contexto emocional, pode representar o "esforço externo necessário" ou a "energia emocional trocada".
\end{remark}

\section{Teorema Fundamental para Integrais de Linha}

\begin{theorem}[Teorema Fundamental - Versão Emocional]
Seja $V: \mathcal{E} \to \mathbb{R}$ uma função diferenciável (custo existencial). Então, para qualquer curva $\gamma$ de $x_0$ a $x_1$:
\[
\int_\gamma \nabla V \cdot d\vec{r} = V(x_1) - V(x_0)
\]
Em particular, a integral depende apenas dos pontos inicial e final, não do caminho.
\end{theorem}

\begin{proof}
Pela regra da cadeia:
\[
\frac{d}{dt} V(\gamma(t)) = \nabla V(\gamma(t)) \cdot \gamma'(t)
\]
Integrando de $a$ a $b$:
\[
\int_a^b \nabla V(\gamma(t)) \cdot \gamma'(t) \, dt = \int_a^b \frac{d}{dt} V(\gamma(t)) \, dt = V(\gamma(b)) - V(\gamma(a))
\]
\end{proof}

\begin{corollary}[Conservação do Custo em Ciclos]
Se $\gamma$ é um caminho fechado ($x_0 = x_1$), então:
\[
\oint_\gamma \nabla V \cdot d\vec{r} = 0
\]
\end{corollary}

\subsection{Campos Conservativos e Potenciais}

\begin{definition}[Campo Conservativo]
Um campo vetorial $F$ é \textbf{conservativo} se existe $V: \mathcal{E} \to \mathbb{R}$ tal que $F = \nabla V$. $V$ é chamado \textbf{potencial} de $F$.
\end{definition}

\begin{theorem}[Equivalências para Campos Conservativos]
Para um campo $F$ em um domínio simplesmente conexo $\mathcal{D} \subseteq \mathcal{E}$, são equivalentes:
\begin{enumerate}
\item $F$ é conservativo ($F = \nabla V$ para algum $V$)
\item $\oint_\gamma F \cdot d\vec{r} = 0$ para toda curva fechada $\gamma$ em $\mathcal{D}$
\item $\int_\gamma F \cdot d\vec{r}$ depende apenas dos pontos final e inicial
\end{enumerate}
\end{theorem}

\begin{example}[Campo Não-Conservativo: Trauma como Singularidade]
Considere o campo:
\[
F(x,y) = \left( \frac{-y}{x^2+y^2}, \frac{x}{x^2+y^2} \right)
\]
Este campo não é conservativo em $\mathbb{R}^2 \setminus \{(0,0)\}$. Se $\gamma$ é um círculo centrado na origem:
\[
\oint_\gamma F \cdot d\vec{r} = 2\pi \neq 0
\]
Interpretação: O ponto $(0,0)$ atua como uma \textbf{singularidade traumática} - um ponto que não pode ser integrado, representando um trauma que "prende" a dinâmica emocional.
\end{example}

\section{Teorema de Green: Circulação e Fluxo}

\subsection{Formulação Clássica}

\begin{theorem}[Teorema de Green]
Seja $R \subset \mathbb{R}^2$ uma região simples com fronteira $\partial R$ orientada positivamente. Para campos $P, Q: R \to \mathbb{R}$ suaves:
\[
\oint_{\partial R} (P\,dx + Q\,dy) = \iint_R \left( \frac{\partial Q}{\partial x} - \frac{\partial P}{\partial y} \right) dA
\]
\end{theorem}

\subsection{Interpretação Emocional}

\begin{center}
\begin{tikzpicture}[scale=1.2]
  % Região R
  \fill[green!20] plot[smooth cycle] coordinates {(1,1) (3,1) (3.5,2) (2.5,3) (1.5,2.5)};
  \draw[thick, green] plot[smooth cycle] coordinates {(1,1) (3,1) (3.5,2) (2.5,3) (1.5,2.5)};
  \node at (2.5,2) {$R$ (região emocional)};
  
  % Fronteira
  \draw[red, thick, ->, decoration={markings, mark=at position 0.2 with {\arrow{>}}, 
        mark=at position 0.5 with {\arrow{>}}, mark=at position 0.8 with {\arrow{>}}}, 
        postaction={decorate}] plot[smooth cycle] coordinates {(1,1) (3,1) (3.5,2) (2.5,3) (1.5,2.5)};
  \node at (1,0.7) {$\partial R$ (fronteira)};
  
  % Rotacional
  \foreach \x in {1.5,2,2.5,3}
    \foreach \y in {1.5,2,2.5}
      \draw[blue, ->] (\x,\y) -- +(0.3,0);
  
  \node at (3.5,2.5) {$\text{rot } F$};
\end{tikzpicture}
\end{center}

\begin{definition}[Circulação e Rotacional]
Para um campo $F = (P, Q)$:
\begin{itemize}
\item \textbf{Circulação}: $\oint_{\partial R} F \cdot d\vec{r}$ mede a tendência do campo a circular em torno de $R$
\item \textbf{Rotacional}: $\text{rot } F = \frac{\partial Q}{\partial x} - \frac{\partial P}{\partial y}$ mede a densidade de circulação
\end{itemize}
\end{definition}

\begin{example}[Ciclos Emocionais]
Considere um campo representando influências sociais:
\[
F(x,y) = (-y, x)
\]
O rotacional é $\text{rot } F = 2$. Para um disco $D$ de raio $R$:
\[
\oint_{\partial D} F \cdot d\vec{r} = \iint_D 2 \, dA = 2\pi R^2
\]
Interpretação: A região tem forte circulação emocional - padrões cíclicos persistentes.
\end{example}

\section{Teorema de Stokes: Generalização para Superfícies}

\subsection{Formulação Geral}

\begin{theorem}[Teorema de Stokes]
Seja $S$ uma superfície orientada com fronteira $\partial S$. Para um campo vetorial $F$ suave:
\[
\oint_{\partial S} F \cdot d\vec{r} = \iint_S (\nabla \times F) \cdot \vec{n} \, dS
\]
onde $\nabla \times F$ é o rotacional e $\vec{n}$ é o vetor normal.
\end{theorem}

\subsection{Aplicação: Buracos na Memória Emocional}

\begin{center}
\begin{tikzpicture}[scale=1.2]
  % Superfície tipo "rosquinha"
  \shade[ball color=blue!20, opacity=0.7] (2,1.5) circle (1.5);
  \draw[thick] (2,1.5) circle (1.5);
  \draw[thick] (2,1.5) circle (0.5);
  
  % Fronteiras
  \draw[red, thick, ->] (3.5,1.5) arc (0:360:1.5 and 0.7);
  \draw[blue, thick, ->] (2.5,1.5) arc (0:-360:0.5 and 0.2);
  
  % Vetores normais
  \foreach \x in {1,2,3}
    \draw[->, green!50!black] ({\x+0.5},{1.5+0.3*sin(\x*60)}) -- ++(0,0.5);
  
  \node at (3.5,2.5) {$S$ (superfície emocional)};
  \node at (3.5,1.5) {$\partial S_1$};
  \node at (2.5,1.2) {$\partial S_2$};
  \node at (0.5,2) {$\vec{n}$};
\end{tikzpicture}
\end{center}

\begin{example}[Trauma como Buraco Topológico]
Considere uma superfície $S$ com um buraco (como uma rosquinha). O campo:
\[
F(x,y,z) = \left( \frac{-y}{x^2+y^2}, \frac{x}{x^2+y^2}, 0 \right)
\]
tem rotacional zero em $S$, mas:
\[
\oint_{\partial S_1} F \cdot d\vec{r} = 2\pi, \quad \oint_{\partial S_2} F \cdot d\vec{r} = -2\pi
\]
Interpretação: O buraco representa uma \textbf{singularidade traumática} que não pode ser integrada. A memória emocional tem uma "lacuna" que perturba a dinâmica.
\end{example}

\section{Teorema da Divergência: Fluxo e Fontes}

\subsection{Formulação Clássica}

\begin{theorem}[Teorema da Divergência (Gauss)]
Seja $V \subset \mathbb{R}^3$ um sólido com fronteira $\partial V$. Para um campo vetorial $F$ suave:
\[
\iint_{\partial V} F \cdot \vec{n} \, dS = \iiint_V (\nabla \cdot F) \, dV
\]
\end{theorem}

\subsection{Interpretação Emocional: Fontes e Sumidouros}

\begin{definition}[Divergência Emocional]
Para um campo $F = (P, Q, R)$:
\[
\nabla \cdot F = \frac{\partial P}{\partial x} + \frac{\partial Q}{\partial y} + \frac{\partial R}{\partial z}
\]
mede a taxa de "expansão" ou "contração" do campo.
\end{definition}

\begin{center}
\begin{tikzpicture}[scale=1.2]
  % Volume
  \draw[thick] (0,0) -- (3,0) -- (3,2) -- (0,2) -- cycle;
  \draw[thick] (0,0) -- (1,1) -- (4,1) -- (3,0);
  \draw[thick] (3,0) -- (4,1) -- (4,3) -- (3,2);
  \draw[thick] (0,2) -- (1,3) -- (4,3) -- (3,2);
  \draw[thick] (1,3) -- (1,1);
  \draw[thick] (4,1) -- (4,3);
  
  % Vetores normais na superfície
  \foreach \x in {0.5,1.5,2.5}
    \foreach \y in {0.5,1.5}
      \draw[->, red] (\x,\y,0) -- ++(0,0,0.7);
  
  % Fonte interna
  \draw[->, blue, thick] (2,1,1.5) -- ++(0.5,0.3,0.2);
  \draw[->, blue, thick] (2,1,1.5) -- ++(-0.3,0.5,0.1);
  \draw[->, blue, thick] (2,1,1.5) -- ++(0.2,-0.4,0.5);
  \node at (2,1,1.5) [circle, fill=blue, inner sep=2pt] {};
  \node at (2.5,1.8,1.5) {$\nabla \cdot F > 0$ (fonte)};
  
  \node at (2, -0.5) {$V$ (volume emocional)};
  \node at (-0.5, 1) {$\partial V$ (fronteira)};
\end{tikzpicture}
\end{center}

\begin{example}[Meditação como Sumidouro]
O campo de meditação $F_{\text{med}}(x) = -\varepsilon x$ tem divergência:
\[
\nabla \cdot F_{\text{med}} = -3\varepsilon < 0
\]
Interpretação: A meditação atua como um \textbf{sumidouro} que reduz a intensidade emocional. O fluxo total para fora de qualquer volume é negativo, indicando dissipação.
\end{example}

\section{Homotopia e Grupos Fundamentais}

\subsection{Trajetórias Homotópicas}

\begin{definition}[Homotopia de Caminhos]
Dois caminhos $\gamma_0, \gamma_1: [0,1] \to \mathcal{E}$ com mesmos pontos inicial e final são \textbf{homotópicos} se existe uma função contínua $H: [0,1] \times [0,1] \to \mathcal{E}$ tal que:
\[
H(s,0) = \gamma_0(s), \quad H(s,1) = \gamma_1(s), \quad H(0,t) = x_0, \quad H(1,t) = x_1
\]
\end{definition}

\begin{center}
\begin{tikzpicture}[scale=1.2]
  % Dois caminhos homotópicos
  \draw[red, thick] plot[smooth, tension=0.7] coordinates {(0,0) (1,1.5) (3,1) (4,0)};
  \draw[blue, thick] plot[smooth, tension=0.7] coordinates {(0,0) (1,0.5) (2.5,2) (4,0)};
  
  % Pontos inicial e final
  \fill[black] (0,0) circle (2pt) node[left] {$x_0$};
  \fill[black] (4,0) circle (2pt) node[right] {$x_1$};
  
  % Homotopia
  \foreach \t in {0.2,0.4,0.6,0.8}
    \draw[green!50!black, dashed, opacity=0.5] 
      plot[smooth, tension=0.7] coordinates {(0,0) ({1+0.5*\t},{0.5+\t}) ({2.5+0.5*\t},{2-0.5*\t}) (4,0)};
  
  \node at (2,2.5) {Família contínua de caminhos};
  \node at (1, -0.7) {$\gamma_0$ e $\gamma_1$ são homotópicos};
\end{tikzpicture}
\end{center}

\begin{theorem}[Invariância por Homotopia]
Se $F$ é um campo conservativo, então:
\[
\int_{\gamma_0} F \cdot d\vec{r} = \int_{\gamma_1} F \cdot d\vec{r}
\]
para quaisquer caminhos homotópicos $\gamma_0, \gamma_1$.
\end{theorem}

\subsection{Grupo Fundamental do Espaço Emocional}

\begin{definition}[Grupo Fundamental]
O \textbf{grupo fundamental} $\pi_1(\mathcal{E}, x_0)$ é o conjunto de classes de homotopia de loops baseados em $x_0$, com a operação de concatenação.
\end{definition}

\begin{example}[Espaços Simplesmente Conexos vs. Não-Simplesmente Conexos]
\begin{itemize}
\item $\mathbb{R}^n$ é simplesmente conexo: $\pi_1(\mathbb{R}^n) = 0$
\item $\mathbb{R}^2 \setminus \{0\}$ não é simplesmente conexo: $\pi_1(\mathbb{R}^2 \setminus \{0\}) = \mathbb{Z}$
\end{itemize}
Interpretação: Um buraco no espaço emocional (trauma) torna o grupo fundamental não-trivial, permitindo diferentes classes de caminhos que não podem ser deformados continuamente uns nos outros.
\end{example}

\section{Aplicações à Dinâmica Emocional}

\subsection{Diagnóstico de Traumas via Integração}

\begin{algorithm}[H]
\caption{Detecção de Singularidades Traumáticas}
\begin{enumerate}
\item Escolha um loop $\gamma$ no espaço emocional
\item Calcule $\oint_\gamma F \cdot d\vec{r}$ para o campo emocional $F$
\item Se $\oint_\gamma F \cdot d\vec{r} \neq 0$:
   \begin{itemize}
   \item Existe uma singularidade dentro de $\gamma$
   \item O valor da integral indica a "força" do trauma
   \item Múltiplos loops podem mapear a topologia do trauma
   \end{itemize}
\item Repita para diferentes loops para localizar singularidades
\end{enumerate}
\end{algorithm}

\subsection{Teorema do Trabalho Mínimo}

\begin{theorem}[Caminho de Menor Custo]
Dados dois estados $x_0, x_1 \in \mathcal{E}$, o caminho que minimiza $\int_\gamma V \, ds$ é o \textbf{caminho geodésico} em relação à métrica:
\[
g_{ij}(x) = V(x) \delta_{ij}
\]
\end{theorem}

\begin{proof}[Ideia]
Minimizar $\int_a^b V(\gamma(t))\|\gamma'(t)\| dt$ é um problema variacional. As equações de Euler-Lagrange dão:
\[
\frac{d}{dt}\left( V(\gamma(t)) \frac{\gamma'(t)}{\|\gamma'(t)\|} \right) = \|\gamma'(t)\| \nabla V(\gamma(t))
\]
\end{proof}

\subsection{Invariantes Topológicos da Memória}

\begin{definition}[Número de Enrolamento]
Para um loop $\gamma$ em torno de uma singularidade $p$:
\[
W(\gamma, p) = \frac{1}{2\pi} \oint_\gamma \frac{(x-p_x)dy - (y-p_y)dx}{(x-p_x)^2 + (y-p_y)^2}
\]
mede quantas vezes $\gamma$ circula $p$.
\end{definition}

\begin{theorem}[Invariância Topológica]
$W(\gamma, p)$ é invariante por homotopia e detecta singularidades.
\end{theorem}

\section{Conclusão: Geometria Integral da Experiência}

As ferramentas de integração de linha e os teoremas fundamentais do cálculo vetorial fornecem uma linguagem poderosa para analisar a dinâmica emocional:

\begin{enumerate}
\item \textbf{Quantificação de Experiências}: Integrais de linha medem custos acumulados e trabalhos emocionais

\item \textbf{Detecção de Estruturas Ocultas}: Singularidades (traumas) manifestam-se como obstruções à integrabilidade

\item \textbf{Classificação Topológica}: Grupos fundamentais classificam espaços emocionais por sua "conectividade emocional"

\item \textbf{Otimização de Trajetórias}: Princípios variacionais identificam caminhos de menor custo

\item \textbf{Conservação e Dissipação}: Teoremas integrais relacionam comportamentos locais e globais
\end{enumerate}

\subsection{Perspectivas Futuras}

\begin{itemize}
\item \textbf{Geometria Riemanniana Emocional}: Métricas que dependem do estado emocional
\item \textbf{Cohomologia de De Rham}: Classificação de campos não-conservativos por classes cohomológicas
\item \textbf{Integrais de Caminho de Feynman}: Quantização da dinâmica emocional
\item \textbf{Topologia Persistente}: Análise de dados emocionais via homologia persistente
\item \textbf{Teoria de Morse}: Análise de pontos críticos em paisagens emocionais
\end{itemize}

\section*{Apêndice: Fórmulas Importantes}

\begin{table}[H]
\centering
\caption{Resumo dos Teoremas Integrais}
\begin{tabular}{p{4cm}p{6cm}p{4cm}}
\toprule
\textbf{Teorema} & \textbf{Fórmula} & \textbf{Interpretação Emocional} \\
\midrule
Fundamental & $\int_\gamma \nabla V \cdot d\vec{r} = V(b) - V(a)$ & Diferença de custo \\
\hline
Green & $\oint_{\partial R} (P dx + Q dy) = \iint_R (Q_x - P_y) dA$ & Circulação vs. rotação \\
\hline
Stokes & $\oint_{\partial S} F \cdot d\vec{r} = \iint_S (\nabla \times F) \cdot \vec{n} dS$ & Circulação em superfícies \\
\hline
Gauss & $\iint_{\partial V} F \cdot \vec{n} dS = \iiint_V (\nabla \cdot F) dV$ & Fluxo vs. fontes \\
\bottomrule
\end{tabular}
\end{table}

\end{document}